\documentclass[aspectratio=169, 11pt]{beamer}

\usepackage[utf8]{inputenc}
\usepackage[T5]{fontenc}
\usepackage[vietnamese]{babel}
\usepackage{amsmath, amssymb, amsfonts}
\usepackage{array, tabularx, booktabs}
\usepackage{xcolor}

\setbeamertemplate{navigation symbols}{}
\setbeamercolor{background canvas}{bg=white}
\setbeamertemplate{itemize item}[circle]

\definecolor{uitblue}{RGB}{0, 0, 204}
\definecolor{footleft}{RGB}{0, 0, 204}
\definecolor{footmid}{RGB}{0, 0, 153}
\definecolor{footright}{RGB}{0, 0, 115}

\setbeamercolor{frametitle}{fg=uitblue}
\setbeamerfont{frametitle}{series=\bfseries, size=\large}

\setbeamercolor{footleftbox}{bg=footleft, fg=white}
\setbeamercolor{footmidbox}{bg=footmid, fg=white}
\setbeamercolor{footrightbox}{bg=footright, fg=white}

\setbeamertemplate{footline}{%
  \leavevmode%
  \hbox{%
    \begin{beamercolorbox}[wd=.333333\paperwidth,ht=2.5ex,dp=1ex,left]{footleftbox}%
      \hspace*{2ex}\textbf{\color{white}\small TS. Nguyễn Văn Hợi}%
    \end{beamercolorbox}%
    \begin{beamercolorbox}[wd=.333333\paperwidth,ht=2.5ex,dp=1ex,center]{footmidbox}%
      \textbf{\color{white}\small UIT -- Cấu Trúc Rời Rạc}%
    \end{beamercolorbox}%
    \begin{beamercolorbox}[wd=.333334\paperwidth,ht=2.5ex,dp=1ex,right]{footrightbox}%
      \textbf{\color{white}\small \insertframenumber{} / 35}\hspace*{2ex}%
    \end{beamercolorbox}%
  }%
  \vskip0pt%
}

\begin{document}

\begin{frame}{Mệnh đề là gì?}
Một \textbf{mệnh đề} là một câu khẳng định một sự kiện: nó phải có giá trị \textbf{Đúng} hoặc \textbf{Sai}, nhưng không thể vừa đúng vừa sai.

\vspace{0.2cm}
\textit{Câu hỏi: Các câu sau có phải là mệnh đề không?}

\vspace{0.2cm}
\textbf{Ví dụ về Mệnh đề:}
\begin{itemize}
  \item ``Hà Nội là thủ đô của Việt Nam.'' $\to$ Mệnh đề Đúng.
  \item ``$0 > e$'' hoặc ``$0 > 2$.'' $\to$ Mệnh đề Sai.
  \item ``$3$ là một số nguyên tố.'' $\to$ Mệnh đề Đúng.
\end{itemize}

\textbf{Không phải Mệnh đề:}
\begin{itemize}
  \item ``Một ngày đẹp trời làm sao!'' (Câu cảm thán).
  \item ``Thầy Hợi rất đẹp trai!'' (Chủ quan/Cảm thán).
  \item ``$3 + x = 2$.'' $\to$ Không phải mệnh đề (Chưa xác định được Đúng/Sai).
\end{itemize}
\end{frame}

\begin{frame}{Ký hiệu và Phân loại}
\begin{itemize}
  \item \textbf{Ký hiệu:} $p, q, r, \dots$
  \item \textbf{Giá trị chân lý:} Đ (Đúng/$1$) hoặc S (Sai/$0$).
\end{itemize}

\textbf{Phân loại:}
\begin{enumerate}
  \item \textbf{Mệnh đề sơ cấp (Đơn giản):} Chỉ chứa một sự kiện duy nhất. Không thể chia nhỏ hơn.
  \begin{itemize}
    \item ``$3$ là một số nguyên tố.''
  \end{itemize}
  \item \textbf{Mệnh đề phức hợp:} Được tạo ra bằng cách kết hợp các mệnh đề sơ cấp bằng các liên từ/toán tử logic (không, và, hoặc, nếu\dots thì, nếu và chỉ nếu).
  \begin{itemize}
    \item ``$2$ là số nguyên tố và là một số chẵn.''
    \item ``Thầy Hợi đang dạy học hoặc đang xem phim.''
  \end{itemize}
\end{enumerate}
\vspace{0.3cm}
Về mặt toán học, mệnh đề phức hợp được hình thành từ các mệnh đề hiện có bằng các toán tử logic.
\end{frame}

\begin{frame}[plain]
  \centering
  \vspace*{2.2cm}
  {\Huge\textbf{\color{uitblue}2. Các toán tử Logic}}
  \par\vspace{0.4cm}
  {\color{uitblue}\rule{0.55\textwidth}{1.2pt}}
\end{frame}

\begin{frame}{1. Phép phủ định (NOT)}
Phép \textbf{phủ định} của $p$ được ký hiệu là $\neg p$ hoặc $\bar{p}$, là mệnh đề: ``Không phải trường hợp $p$''. Giá trị chân lý của nó ngược lại với $p$.

\vspace{0.4cm}
\begin{columns}[c, onlytextwidth]
  \begin{column}{0.35\textwidth}
    \centering
    \begin{tabular}{|c|c|}
    \hline
    $p$ & $\bar{p}$ \\
    \hline
    0 & 1 \\
    1 & 0 \\
    \hline
    \end{tabular}
  \end{column}
  \begin{column}{0.65\textwidth}
    \textbf{Ví dụ:}
    \begin{itemize}
      \item $p$: ``Paris là thủ đô của Pháp.''\\
            $\to \bar{p}$: ``Paris \textbf{không phải} là thủ đô của Pháp.''
      \item $q$: ``$9$ là một số chính phương.''\\
            $\to \bar{q}$: ``$9$ \textbf{không phải} là một số chính phương.''
    \end{itemize}
  \end{column}
\end{columns}
\end{frame}

\begin{frame}{2. Phép hoặc (OR) và 3. Phép giao (AND)}
\begin{columns}[t, onlytextwidth]
  \begin{column}{0.51\textwidth}
    \textbf{Phép tuyển} ($p \lor q$) là mệnh đề: ``$p$ hoặc $q$''.\\
    \textit{``Đúng nếu có ít nhất một mệnh đề đúng.''}

    \vspace{0.3cm}
    \centering
    \begin{tabular}{|c|c|c|}
    \hline
    $p$ & $q$ & $p \lor q$ \\
    \hline
    0 & 0 & 0 \\
    0 & 1 & 1 \\
    1 & 0 & 1 \\
    1 & 1 & 1 \\
    \hline
    \end{tabular}

    \vspace{0.3cm}
    \raggedright
    \textit{VD: ``$3$ là số nguyên tố \textbf{hoặc} $3$ là số chẵn.'' (Đúng)}
  \end{column}
  \begin{column}{0.47\textwidth}
    \textbf{Phép hội} ($p \land q$) là mệnh đề: ``$p$ và $q$''.\\
    \textit{``Đúng chỉ khi cả hai đều đúng.''}

    \vspace{0.3cm}
    \centering
    \begin{tabular}{|c|c|c|}
    \hline
    $p$ & $q$ & $p \land q$ \\
    \hline
    0 & 0 & 0 \\
    0 & 1 & 0 \\
    1 & 0 & 0 \\
    1 & 1 & 1 \\
    \hline
    \end{tabular}

    \vspace{0.3cm}
    \raggedright
    \textit{VD: ``Trời đã mưa \textbf{và} hôm đó là thứ Sáu.''}
  \end{column}
\end{columns}
\end{frame}

\begin{frame}{4. Phép tuyển loại (XOR)}
Phép \textbf{tuyển loại} của $p$ và $q$, ký hiệu là $p \oplus q$, là mệnh đề đúng khi \textbf{chỉ duy nhất một} trong hai mệnh đề $p$ hoặc $q$ đúng, và sai trong các trường hợp còn lại.

\vspace{0.3cm}
\begin{columns}[t, onlytextwidth]
  \begin{column}{0.4\textwidth}
    \textbf{Bảng chân lý:}
    \vspace{0.2cm}
    \begin{center}
    \begin{tabular}{|c|c|c|}
    \hline
    $p$ & $q$ & $p \oplus q$ \\
    \hline
    0 & 0 & 0 \\
    0 & 1 & 1 \\
    1 & 0 & 1 \\
    1 & 1 & 0 \\
    \hline
    \end{tabular}
    \end{center}
    \vspace{0.1cm}
    \centering
    \textit{Khác với OR ở dòng cuối!}
  \end{column}
  \begin{column}{0.58\textwidth}
    \textbf{Ví dụ:} \textit{``Sinh viên có thể chọn súp miễn phí hoặc salad miễn phí kèm bữa ăn.''}
    \begin{itemize}
      \item $p$: ``Bạn chọn súp.''
      \item $q$: ``Bạn chọn salad.''
    \end{itemize}
    \textbf{Logic ($p \oplus q$):}
    \begin{itemize}
      \item Đúng nếu bạn chỉ lấy một thứ.
      \item \textbf{Sai} nếu bạn cố lấy \textbf{cả hai} hoặc không lấy thứ nào.
    \end{itemize}
  \end{column}
\end{columns}
\end{frame}

\begin{frame}{5. Phép kéo theo và 6. Phép tương đương}
\begin{columns}[t, onlytextwidth]
  \begin{column}{0.48\textwidth}
    \textbf{Phép kéo theo} ($p \rightarrow q$):\\
    ``Nếu $p$, thì $q$.''

    \vspace{0.25cm}
    \centering
    \begin{tabular}{|c|c|c|}
    \hline
    $p$ & $q$ & $p \rightarrow q$ \\
    \hline
    0 & 0 & 1 \\
    0 & 1 & 1 \\
    1 & 0 & 0 \\
    1 & 1 & 1 \\
    \hline
    \end{tabular}
  \end{column}
  \begin{column}{0.48\textwidth}
    \textbf{Phép tương đương} ($p \leftrightarrow q$):\\
    ``$p$ nếu và chỉ nếu $q$.''

    \vspace{0.25cm}
    \centering
    \begin{tabular}{|c|c|c|}
    \hline
    $p$ & $q$ & $p \leftrightarrow q$ \\
    \hline
    0 & 0 & 1 \\
    0 & 1 & 0 \\
    1 & 0 & 0 \\
    1 & 1 & 1 \\
    \hline
    \end{tabular}
  \end{column}
\end{columns}

\vspace{0.4cm}
\textit{Ví dụ $p \rightarrow q$: ``Nếu bạn học chăm chỉ, bạn sẽ qua môn.''}\\
\textit{Ví dụ $p \leftrightarrow q$: ``Bạn có thể đi chuyến bay nếu và chỉ nếu bạn mua vé.''}
\end{frame}

\begin{frame}{Hiểu phép kéo theo qua một lời hứa}
Thầy hứa: ``\textit{Nếu em học chăm chỉ} ($p$) \textit{thì em sẽ qua môn} ($q$).''\\
\textbf{Câu hỏi:} Trong trường hợp nào thầy đã \textbf{nói sai}?

\vspace{0.25cm}
\begin{center}
\begin{tabular}{|c|c|l|c|}
\hline
$p$ & $q$ & \textbf{Tình huống} & $p \rightarrow q$ \\
\hline
1 & 1 & Học chăm, qua môn $\Rightarrow$ thầy giữ lời & 1 \\
1 & 0 & Học chăm, vẫn trượt $\Rightarrow$ \textcolor{red}{thầy thất hứa} & \textcolor{red}{0} \\
0 & 1 & Không học, vẫn qua $\Rightarrow$ thầy không nói sai & 1 \\
0 & 0 & Không học, trượt $\Rightarrow$ thầy không nói sai & 1 \\
\hline
\end{tabular}
\end{center}

\vspace{0.25cm}
\textbf{\textcolor{red}{Ghi nhớ:}}\\
$p \rightarrow q$ chỉ \textbf{Sai} khi $p$ \textbf{Đúng} mà $q$ \textbf{Sai}. Giả thiết sai thì phép kéo theo luôn đúng.
\end{frame}

\begin{frame}{Kiểm tra nhanh}
Cho $p$ là \textbf{Đúng} ($1$) và $q$ là \textbf{Sai} ($0$). Tính giá trị chân lý:

\vspace{0.4cm}
\begin{columns}[T, onlytextwidth]
  \begin{column}{0.48\textwidth}
    \centering
    \begin{tabular}[t]{lll}
      \textcolor{uitblue}{\textbf{1.}} & $\neg p$ & $= \textcolor{red}{0}$ \\[6pt]
      \textcolor{uitblue}{\textbf{2.}} & $p \lor q$ & $= \textcolor{green!60!black}{1}$ \\[6pt]
      \textcolor{uitblue}{\textbf{3.}} & $p \land q$ & $= \textcolor{red}{0}$ \\[6pt]
      \textcolor{uitblue}{\textbf{4.}} & $p \oplus q$ & $= \textcolor{green!60!black}{1}$
    \end{tabular}
  \end{column}
  \begin{column}{0.48\textwidth}
    \centering
    \begin{tabular}[t]{lll}
      \textcolor{uitblue}{\textbf{5.}} & $p \rightarrow q$ & $= \textcolor{red}{0}$ \\[6pt]
      \textcolor{uitblue}{\textbf{6.}} & $q \rightarrow p$ & $= \textcolor{green!60!black}{1}$ \\[6pt]
      \textcolor{uitblue}{\textbf{7.}} & $p \leftrightarrow q$ & $= \textcolor{red}{0}$
    \end{tabular}
  \end{column}
\end{columns}

\vspace{0.4cm}
\centering
\textbf{Nhận xét:} $p \rightarrow q$ và $q \rightarrow p$ không tương đương!
\end{frame}

\end{document}
